% !PRINTED 774
\weberpage{774}
Nun setzen wir, indem wir unter $t$ eine Variable verstehen:

\[
\begin{aligned}
{}[t - \chi(1)][t - \chi(r)] \ldots [t - \chi(r^{n-1})] &= \\
t^n + A_1t^{n-1} + A_2t^{n-2} + \cdots + A_n, &
\end{aligned}
\]

worin die Coëfficienten $A_1, A_2, \ldots, A_n$ als ganze symmetrische Functionen der $n$ Grössen $\chi(1), \chi(r), \ldots, \chi(r^{n-1})$ ganze rationale Zahlen sind. Wenn wir aber darin

\[
t = \chi(r) = \dfrac{[A(\alpha)]^n}{p}
\]

setzen, so verschwindet die linke Seite, und es ergiebt sich durch Multiplication mit $p^n$:

\weq{14}{[A(\alpha)]^{n^2} = p\,C(\alpha),}

wenn $C(\alpha)$ eine Zahl von der Form (9) mit \gesperrt{ganzzahligen Coëfficienten} bedeutet.

Wenn man die Formel (14) wiederholt in die $p^{\text{te}}$ Potenz erhebt und den Satz (1) anwendet, so ergiebt sich daraus für jede Potenz, die nicht kleiner als $n^2$ ist:

\weq{15}{A(\alpha^{p^\lambda}) = p\,\xi(\alpha),}

worin die ganze Function $\xi(\alpha)$ gleichfalls ganzzahlige Coëfficienten hat.

Hierin können wir $\lambda$ durch $\varphi(m)$ theilbar annehmen, und dann ist, da $m$ durch $p$ nicht theilbar ist, nach dem verallgemeinerten \gesperrt{Fermat}'schen Satze (Bd.~II, \S.\,14):

\[
p^\lambda \equiv 1 \ (\mathrm{mod}\ m), \quad \alpha^{p^\lambda} = \alpha,
\]

so dass aus (14) folgt: % [sic] im Druck; die Gleichung (16) folgt sachlich aus (15), nicht aus (14)

\weq{16}{A(\alpha) = p\,\xi(\alpha).}

Es sind daher in (13) die Coëfficienten $a_0, a_1, \ldots, a_{\nu-1}$ alle durch $p$ theilbar und $a$ ist folglich nicht durch $p$ theilbar. Ebenso wird gezeigt, dass $b_0, b_1, \ldots, b_{\nu-1}$ durch $p$ theilbar sind, während $b$ durch $p$ untheilbar ist.

Demnach ergiebt sich aus (12), wenn $A(\alpha)B(\alpha) = p^2\Theta(\alpha)$ gesetzt wird:

\weq{17}{ab = p\,\Theta(\alpha),}

worin $\Theta(\alpha)$ eine Function von $\alpha$ von der Form (9) mit ganzzahligen Coëfficienten bedeutet. Die Gleichung (17) müsste aber, da sie höchstens vom $(\nu-1)^{\text{ten}}$ Grade ist, in Bezug auf $\alpha$ identisch sein, und es müsste $ab$ durch $p$ theilbar sein, was nicht möglich ist.
